\chapter{Matériel}
Les données sélectionnées servent de proxy pour les variables environnementales clés dans la structuration des écosystèmes. L'hydrodynamisme est approximé par le Finite-Time Lyapunov Exponent, tandis que la température et la salinité rendent compte des conditions physico-chimiques du milieu. L'activité phytoplanctonique est estimée à partir des concentrations de neuf pigments photosynthétiques et la distribution des niveaux trophiques intermédiaires est évaluée par acoustique active.

\section{Zone d'étude et plan d'échantillonnage}

L'étude a été réalisée sur les années 2018, 2021, 2022 et 2023, en saison estivale (janvier-mars) afin de limiter la variabilité saisonnière, et de jour (calculé selon l'angle solaire), afin de limiter l'effet de la migration verticale journalière (DVM) \footnote{Déplacement synchronisé dans la colonne d'eau au rythme du cycle jour/nuit : les organismes résident en profondeur durant le jour et remontent vers la surface la nuit, principalement pour se nourrir tout en réduisant le risque de prédation visuelle en eaux éclairées. Ce phénomène est considéré comme le plus grand mouvement synchronisé de biomasse à l'échelle de la planète. \cite{Bandara2021}}. La fenêtre d'étude a été fixée à [45,90]°E de longitude et [30, 60]°S de latitude afin de cibler précisémment le secteur indien de l'Océan Austral.

\section{Données d'acoustique active (OBSAUSTRAL)}

Les données d'acoustiques actives ont été collectées par transects \footnote{Le bateau se déplace et collecte en même temps des données, par opposition aux données "stations" collectées lorsque le bateau est en arrêt). Celà a un coût en terme de résolution des spectrogramme et d'interprétabilité mais permet d'obtenir une plus grande couverture spatiale.} lors des campagnes océanographiques OBSAUSTRAL/THEMISTO 2018, 2021, 2022, 2023 en saison estivale (voir figure \ref{fig:transectsparannee2})\cite{flotteoceanographique_campagnes}.  
\begin{figure}[H]
	\centering
	\includegraphics[width=0.6\textwidth]{images/02_materiel/transects_par_annee_2}
	\caption{Carte des transects effectués par le Marion Dufresnes II lors des campagnes océanographiques OBSAUSTRAL/THEMISTO 2018, 2021, 2022 et 2023.}
	\label{fig:transectsparannee2}
\end{figure}

\vspace{1cm}

Deux fréquences d'émissions ont été sélectionnées: 38kHz, afin de comparer nos résultats avec la littérature \footnote{La méthode d'acoustique active ayant été développée par les pêcheries, la fréquence d'échosondeur de référence est 38kHz (fréquence à laquelle sont les poissons résonnent).} et 120kHz, fréquence à laquelle les Euphausiacés (krill) et Copépodes sont le plus détectable (voir Annexe \ref{annexe:rep_freq}, figure \ref{fig:freqresponse}). Les échogrammes (figures \ref{fig:echogramme38khz2023-01-28} et \ref{fig:echogramme120khz2023-01-28}) donnent un aperçu de la nature des données collectées et de la manière dont l'information contenue par le signal reçu (résolution, organismes visibles) est déterminée par la fréquence de l'onde émise. Une forte intensité de l'echo reçu traduis la présence d'une importante biomasse d'organisme à ce temps, localisation géographique et profondeur.

\begin{figure}[H]
	\centering
	\begin{subfigure}[t]{0.7\textwidth}
		\centering
		\includegraphics[width=\textwidth]{images/02_materiel/echogramme_38kHz_2023-01-28}
		\caption{Fréquence d'émission de 38 kHz}
		\label{fig:echogramme38khz2023-01-28}
	\end{subfigure}
	\vfill
	\begin{subfigure}[t]{0.7\textwidth}
		\centering
		\includegraphics[width=\textwidth]{images/02_materiel/echogramme_120kHz_2023-01-28}
		\caption{Fréquence d'émission de 120 kHz}
		\label{fig:echogramme120khz2023-01-28}
	\end{subfigure}
	\caption{Exemple d'échogramme obtenu le 2023-01-28 entre 18h et 00h aux fréquences (a) 38 kHz et (b) 120 kHz}
	\label{fig:echogramme2023-01-28}
\end{figure}

\section{Concentrations de pigments photosynthétiques (PIGMeANN)}
L'analyse de la couleur des océans à partir d'images satellite permet de d'étudier le phytoplancton sur une échelle spatiale et temporelle large. Des méthodes de machine learning et, plus récemment de deep learning, ont permis la caractérisation de "Phytoplankton Functional Types" et de "Phytoplakton Size Classes"  depuis des images satellites \cite{Alvain2005}, \cite{Li2023DLPPCE}, \cite{Waga2026CSD}.  Les spectres de reflectance des océans dépendent des pigments photosynthétiques qu'ils contiennent. L'idée est d'identifier une empreinte spectrale et de l'associer (par deep learning) avec des mesures \textit{in-situ} de High Performance Liquid Chromatography \footnote{L’échantillon est injecté dans une colonne traversée par un solvant sous haute pression, où ses différents composés sont séparés puis détectés individuellement. Permet de séparer et quantifier les pigments afin d’identifier les différents groupes présents.} et de cytométrie en flux\footnote{Les cellules sont entraînées individuellement devant un laser qui mesure leur diffusion (taille/forme des cellules) et leur fluorescence. On peut ensuite identifier les taxons.} qui caractérisent plus précisemment la présence et structure des communauté phytoplanctoniques. 
\begin{table}[H]
	\centering
	\renewcommand{\arraystretch}{1.3}
	\begin{tabular}{
			>{\raggedright\arraybackslash}p{3.4cm}
			>{\raggedright\arraybackslash}p{2.0cm}
			>{\raggedright\arraybackslash}p{4.8cm}
			>{\raggedright\arraybackslash}p{2.6cm}
		}
		\hline
		\textbf{Pigment} & \textbf{Abréviation} & \textbf{Groupe phytoplanctonique associé (PFD)} & \textbf{Classe de taille (PSC)} \\
		\hline
		Divinyl chlorophylle $a$ & dvChla & Procaryotes picoplanctoniques (\textit{Prochlorococcus}) & Picophytoplancton \\
		Chlorophylle $a$ & Chla & Marqueur de biomasse totale (tous groupes) & -- \\
		Chlorophylle $b$ & Chlb & Chlorophycées / Prasinophycées & Picophytoplancton \\
		Zéaxanthine & Zea & Cyanobactéries (\textit{Synechococcus}) & Picophytoplancton \\
		Fucoxanthine & Fuco & Diatomées (dominantes) & Microphytoplancton \\
		19'-hexanoyloxyfucoxanthine & Hex & Prymnésiophycées (haptophytes) & Nanophytoplancton \\
		19'-butanoyloxyfucoxanthine & But & Prymnésiophycées / pélagophytes & Nanophytoplancton \\
		Alloxanthine & Allo & Cryptophycées & Nanophytoplancton \\
		Péridinine & Per & Dinoflagellés (péridiniens) & Microphytoplancton \\
		\hline
	\end{tabular}
	\renewcommand{\arraystretch}{1}
	\caption{Pigments photosynthétiques analysés et groupes phytoplanctoniques hypothétiquement associés (PFD et PSC) dans le secteur de Kerguelen en été austral, d'après \cite{Lasbleiz2014}, \cite{sari2025influence}}
	\label{tab:pigments}
\end{table}

La méthode de deep-learning PIGMeANN a récemment été développée \cite{sari2026regional} : elle permet d'approximer les concentrations de surface de neuf pigments photosynthétiques (tableau \ref{tab:pigments}) sur l'ensemble de la couverture spatiale de l'image satellite donnée en entrée. La figure  \ref{fig:mapchla2023-01-26} illustre la prédiction par PIGMeANN de la Chorophylle-a le 2023-01-26 sur le secteur étudié.

\begin{figure}[H]
	\centering
	\includegraphics[width=0.7\textwidth]{images/02_materiel/map_chla_2023-01-26}
	\caption{Carte de prédiction de la concentration de Chlorophylle-a le 2023-01-26 sur le secteur indient du SO, prédictions issue de l'alorithme PIGMeANN développé par Sari el Dine \cite{sari2026regional}. La présence de nuages gène grandement la prédiction sur la zone totale.}
	\label{fig:mapchla2023-01-26}
\end{figure}

A partir de la connaissance des différentes communautés phytoplanctoniques présente dans la région d'étude, on peut hypothétiquement associer certains pigments à des groupes phytoplanctoniques (PFD) ou Phytoplankton Size Classes (PSC) (voir tableau \ref{tab:pigments})\cite{sari2025influence}, \cite{heidemann2024drivers} \cite{Georges2014}. Ces associations permettent alors, à partir des assemblage de concentration de pigments photosynthétiques prédits par PIGMeANN, de caractériser les structures des communautés phytoplanctoniques en identifiant les PFD et PSC dominantes. 

La chlorophylle-a, utilisée comme proxy de la biomasse phytoplanctonique totale, est exclue de cette identification car peu discriminante : présente à forte concentration chez la quasi-totalité des organismes, elle n'apporte pas d'information sur la composition de la communauté. Pour chaque pigment restant, on calcule donc sa contribution relative (ratio par rapport à la somme des pigments) plutôt que sa concentration absolue :

\begin{equation}
	R_i = \frac{[\text{Pigment}_i]}{\displaystyle\sum_{j \neq \text{Chla}} [\text{Pigment}_j]}
	\label{eq:ratio_pigment}
\end{equation}

\section{Données hydrologiques de température et salinité (CMEMS)}
Les données de températures et salinité sont issues de la plateforme Copernicus \footnote{Programme d’observation de la Terre de l’Union européenne, qui fournit notamment des données et services permettant de surveiller et comprendre l’état de l’environnement, dont l’océan \cite{copernicus2026services}}. Ce sont des données issues des observations temps-réelles CMEMS \footnote{Copernicus Marine Environment Monitoring Service, service marin de Copernicus} et réanalysées par le modèle NEMO \footnote{Nucleus for European Modelling of the Ocean, modèle numérique de l’océan qui simule l’évolution de paramètres comme la température, la salinité, les courants et le niveau de la mer}, et ERA5 \footnote{Modèle atmosphérique de l'ECMWF (European Centre for Medium-Range Weather Forecasts) qui fournit les conditions atmosphériques nécessaires pour forcer le modèle océanique NEMO à sa surface}. Elles ont une résolution spatiale de 0.083° et une résolution temporelle journalière.
Sont téléchargés tous les profils de température et de salinité de la zone d'étude (entre 45-90°E de longitude et 30 et 60°S de latitude), sur l'ensemble de la période étudiée (2018, 2021, 2022, 2023, entre janvier et mars). 

\section{Données hydrodynamiques : Finite Time Lyapunov Exponent}
Le Finite Time Lyapunov Exponent est un indicateur qui, à partir de champs de courants estimés par modélisation, permet de caractériser la séparation ou la convergence de trajectoires de particules d’eau sur une durée finie. On peut ainsi identifier les structures dynamiques susceptibles de contrôler le transport et la dispersion dans l’océan \cite{shadden2005definition}, \cite{dovidio2004mixing}. Son calcul est détaillé en Annexe \ref{annexe:compute_ftle}.

\vspace{0.5em}

\begin{figure}[H]
	\centering
	\includegraphics[width=0.5\textwidth]{images/02_materiel/map_ftle_2023-01-26}
	\caption{Advection Lagrangienne FTLE backward (90j) du 2023-01-26. Les valeurs élevées de FTLE délimitent des crêtes correspondant  à des frontières de transport de l'écoulement (de fronts ou de tourbillons) le long desquelles convergent les masses d'eau (structures cohérentes lagrangiennes attractives)}
	\label{fig:mapftle2023-01-26}
\end{figure}

Dans notre étude, la durée d'advection est de trente jours en backward, c'est-à-dire que les trajectoires sont intégrées vers l'arrière dans le temps (de $t$ à $t - T$). La résolution temporelle est journalière et la résolution spatiale est de 0.05°.On obtient des cartes de FTLE $\in \mathbb{R}^{time \times lon \times lat}$ pour la période (2018, 2021, 2022, 2023 entre janvier et mars) et région étudiée ([45, 90]°E, [30, 60]°S). 

La figure \ref{fig:mapftle2023-01-26} présente une carte de FTLE issue de la journé 2023-01-26. Une valeur élevée du FTLE indique une forte séparation des trajectoires voisines sur la période considérée. Les structures organisées apparaissant comme des crêtes de FTLE peuvent ainsi être utilisées pour identifier des structures physiques réelles dans l'écoulement océanique, appelées Structures Cohérentes Lagrangiennes \cite{shadden2005definition}. 